\begin{manuallistedformatdiagram}{Floating-Point Register Encodings}[fig:floating-point-register-encodings]
\manualformatrowrange{D / Fn[63:0]}{63}{0}{%
\manualformatfield{s}{1}
\manualformatfield{exponent}{11}
\manualformatfield{fraction}{52}
}
\manualformatrowrange{S / Fn[63:0]}{63}{0}{%
\manualformatfield{zero}{32}
\manualformatfield{s}{1}
\manualformatfield{exponent}{8}
\manualformatfield{fraction}{23}
}
\end{manuallistedformatdiagram}
{\small
For a NaN, the most significant fraction bit is the quiet bit: bit 22 for S and bit 51 for D.
\texttt{s} denotes the sign bit. An S result occupies Fn bits 31..0 and writes zero to Fn bits 63..32.
\par}
